\chapter{The Schwinger Model: Two-Dimensional Quantum Electrodynamics}
\label{ch:schwinger-model}
\ConventionsLorentzian

\section{Definition, Gauge Group, and Observable Sector}

The Schwinger model is quantum electrodynamics in \(1+1\) spacetime
dimensions with one massless Dirac field.  We use coordinates
\(x^\mu=(t,x)\),
\[
  \eta_{\mu\nu}=\operatorname{diag}(-,+),
  \qquad
  \epsilon^{01}=+1,
  \qquad
  \gamma=\gamma^0\gamma^1 .
\]
The two-dimensional spinor conventions are those of
Appendix~\ref{app:spinor-conventions}.  Let \(a_\mu\) be a real Abelian
gauge representative and let \(e>0\) be the charge of the dynamical Dirac
field.  The compact \(U(1)\) connection is \(e\,a_\mu\); this convention is
important because external probes will be allowed to have charge \(Q\) with
respect to \(a_\mu\).  Thus
\[
  D_\mu\psi=\partial_\mu\psi-\ii e a_\mu\psi,
  \qquad
  F_{\mu\nu}=\partial_\mu a_\nu-\partial_\nu a_\mu,
  \qquad
  E=F_{01}.
\]
The local representative Lagrangian is
\begin{equation}
  \mathcal L_{\rm Sch}
  =
  -\frac14F_{\mu\nu}F^{\mu\nu}
  -\bar\psi\gamma^\mu D_\mu\psi
  =
  \frac12E^2-\bar\psi\gamma^\mu\partial_\mu\psi
  +e\,a_\mu J^\mu ,
  \label{eq:schwinger-lagrangian}
\end{equation}
where
\begin{equation}
  J^\mu=\ii\bar\psi\gamma^\mu\psi,
  \qquad
  J_A^\mu=\ii\bar\psi\gamma^\mu\gamma\psi .
  \label{eq:schwinger-vector-axial-currents}
\end{equation}
With these definitions \(eJ^\mu\) is the dynamical electric current.  Gauge
representatives are related by
\[
  a_\mu\mapsto a_\mu+\partial_\mu\xi,
  \qquad
  \psi\mapsto \ee^{\ii e\xi}\psi,
  \qquad
  \bar\psi\mapsto\bar\psi\,\ee^{-\ii e\xi}.
\]
On a spatial circle of circumference \(L\), the gauge group includes large
transformations with \(\xi(x+L)-\xi(x)=2\pi n/e\).  Equivalently, the
holonomy \(\exp(\ii e\oint a_x\dd x)\) is the compact gauge-invariant
zero-mode coordinate.  The local computations below may be performed on the
line and then lifted to the circle by including the compact zero mode and
the large-gauge identifications.

A possible theta density is
\begin{equation}
  \mathcal L_\theta=\frac{e\theta}{2\pi}E .
  \label{eq:schwinger-theta-density}
\end{equation}
On a closed Euclidean two-manifold the curvature of \(e\,a\) has integral
periods, so \(\theta\sim\theta+2\pi\).  In the massless theory an anomalous
axial rotation shifts \(\theta\), and the local spectrum is independent of
\(\theta\).  In the massive deformation this shift is no longer a symmetry,
and \(\theta\) becomes a physical parameter.

\begin{definition}[Probe-charge convention]
  A static external probe of charge \(Q\) is defined by the conserved
  external current \(j^\mu_{\rm ext}\) coupled through
  \[
    S_{\rm probe}=\int a_\mu j^\mu_{\rm ext}\,\dd^2x .
  \]
  A straight worldline carries the phase
  \(\exp(\ii Q\int a_\mu\dd x^\mu)\).  The dynamical Dirac field has
  \(Q=e\).  A probe with \(Q/e\in\mathbb Z\) can be screened by dynamical
  charged particles; a probe with nonintegral \(Q/e\) is a fractional probe
  relative to the dynamical charge lattice.
\end{definition}

The canonical momentum conjugate to \(a_x\) is \(E\), while \(a_0\) is a
Lagrange multiplier.  Including an external charge density
\(\rho_{\rm ext}=j^0_{\rm ext}\), variation of \(a_0\) gives the operator
Gauss law
\begin{equation}
  \partial_xE+eJ^0+\rho_{\rm ext}=0 .
  \label{eq:schwinger-gauss-law}
\end{equation}
This equation is a constraint on states and on correlation functions with
line insertions.  It is not optional data added after quantization; it is the
local gauge-invariance condition of the theory.

\begin{definition}[Local sector]
  The local gauge-invariant sector of the massless Schwinger model is the
  local operator algebra generated, after renormalization, by \(F_{\mu\nu}\),
  by gauge-invariant currents such as \(J^\mu\), and by neutral composite
  fields.  The bare fermion field \(\psi\) is not a gauge-invariant local
  observable.  Charged sectors are described by line-dressed charged
  insertions or by states on a spatial manifold with prescribed electric
  flux at infinity.  The exact solution below is first a solution of the
  local gauge-invariant sector and of specified external line probes.
\end{definition}

\section{Current Algebra, Anomaly, and the Electric-Field Equation}

In two dimensions the vector and axial currents are Hodge dual.  With the
gamma-matrix and orientation conventions fixed above,
\begin{equation}
  J_A^\mu=-\epsilon^{\mu\nu}J_\nu .
  \label{eq:schwinger-current-duality}
\end{equation}
The vector current is exactly conserved as a consequence of gauge symmetry:
\[
  \partial_\mu J^\mu=0 .
\]
The axial current is not conserved.  The two-dimensional anomaly formula of
Chapter~\ref{ch:volume-ii-20}, applied to the background connection
\(A_\mu=e a_\mu\), gives the renormalized operator identity
\begin{equation}
  \partial_\mu J_A^\mu
  =
  \frac{e}{2\pi}\epsilon^{\mu\nu}F_{\mu\nu}
  =
  \frac{e}{\pi}E .
  \label{eq:schwinger-axial-anomaly}
\end{equation}
The identity includes the finite contact term chosen in the definition of
the axial-vector current correlator.  It is not the Euler-Lagrange equation
of the classical massless Dirac action.

\begin{proposition}[Anomaly-induced massive electric field]
  In the absence of external line insertions, the gauge-invariant electric
  field of the massless Schwinger model obeys
  \begin{equation}
    (\Box-m_{\rm Sch}^2)E=0,
    \qquad
    m_{\rm Sch}^2=\frac{e^2}{\pi},
    \qquad
    \Box=\partial_\mu\partial^\mu .
    \label{eq:schwinger-electric-field-mass-equation}
  \end{equation}
\end{proposition}

\begin{proof}
  The Maxwell equation obtained from \eqref{eq:schwinger-lagrangian} is
  \[
    \partial_\nu F^{\nu\mu}+eJ^\mu=0 .
  \]
  Since \(F^{10}=E\) and \(F^{01}=-E\), its two components are
  \[
    \partial_xE+eJ^0=0,
    \qquad
    -\partial_tE+eJ^1=0 .
  \]
  Hence
  \[
    J^0=-\frac1e\partial_xE,
    \qquad
    J^1=\frac1e\partial_tE .
  \]
  Using \eqref{eq:schwinger-current-duality},
  \(J_A^0=-J^1\) and \(J_A^1=-J^0\).  Therefore
  \[
    \partial_\mu J_A^\mu
    =
    -\partial_tJ^1-\partial_xJ^0
    =
    \frac1e(-\partial_t^2+\partial_x^2)E
    =
    \frac1e\Box E .
  \]
  Comparing this consequence of Maxwell's equation with the anomaly
  \eqref{eq:schwinger-axial-anomaly} gives
  \(\Box E/e=eE/\pi\), which is \eqref{eq:schwinger-electric-field-mass-equation}.
\end{proof}

This derivation uses only gauge invariance, the exact current anomaly, and
the Maxwell equation.  It already proves that the local electric field
creates a massive neutral particle.  Bosonization gives a stronger statement:
it identifies the entire current-electric-field sector as a free massive
scalar theory and gives exact formulas for line-probe response.

\section{Bosonization of the Current Sector}

The conservation law for \(J^\mu\) implies locally that \(J^\mu\) is the
Hodge dual of a derivative.  The normalization is not arbitrary; it is fixed
by the current algebra of one Dirac fermion.

\begin{theorem}[Current-sector bosonization]
  Let the free massless Dirac field be quantized in the vacuum sector with
  the current normalization \eqref{eq:schwinger-vector-axial-currents}.  The
  gauge-invariant Abelian current subalgebra is represented by a compact
  scalar field \(\varphi\) of period \(\sqrt\pi\), normalized by
  \[
    \mathcal L_\varphi
    =
    -\frac12\partial_\mu\varphi\,\partial^\mu\varphi,
    \qquad
    [\widehat\varphi(t,x),\partial_0\widehat\varphi(t,y)]
    =
    \ii\delta(x-y),
  \]
  through
  \begin{equation}
    J^\mu=\frac1{\sqrt\pi}\epsilon^{\mu\nu}\partial_\nu\varphi,
    \qquad
    J_A^\mu=-\frac1{\sqrt\pi}\partial^\mu\varphi .
    \label{eq:schwinger-bosonization-currents}
  \end{equation}
  Equivalently, after the same local contact-term convention is fixed on
  both sides, the current-generating functional satisfies
  \begin{equation}
    Z_{\rm Dirac}[b]
    =
    \int\mathcal D\varphi\,
    \exp\!\left\{
      \ii\int\dd^2x
      \left[
        -\frac12\partial_\mu\varphi\partial^\mu\varphi
        +\frac1{\sqrt\pi}b_\mu\epsilon^{\mu\nu}\partial_\nu\varphi
      \right]
    \right\}.
    \label{eq:schwinger-current-generating-functional}
  \end{equation}
\end{theorem}

\begin{proof}
  We give the proof for the current sector, which is the part used in the
  gauge theory.  The equal-time commutator of the Dirac current, defined by
  gauge-invariant point splitting and with the vector Ward identity
  preserved, is
  \begin{equation}
    [\widehat J^0(t,x),\widehat J^1(t,y)]
    =
    -\frac{\ii}{\pi}\partial_x\delta(x-y).
    \label{eq:schwinger-current-schwinger-term}
  \end{equation}
  This Schwinger term follows directly from the short-distance singularity
  of the two chiral components of a free massless Dirac field.  If
  \eqref{eq:schwinger-bosonization-currents} holds, then
  \[
    J^0=\frac1{\sqrt\pi}\partial_x\varphi,
    \qquad
    J^1=-\frac1{\sqrt\pi}\partial_t\varphi ,
  \]
  and the canonical scalar commutator gives exactly
  \eqref{eq:schwinger-current-schwinger-term}.  Hence the coefficient in
  \eqref{eq:schwinger-bosonization-currents} is forced.

  It remains to see that this representation gives all current correlators,
  not only their equal-time central term.  Work first in Euclidean signature
  and couple a real background one-form \(b_\mu\) to the conserved vector
  current.  Gauge invariance of the source problem implies that the
  generating functional depends only on the transverse part of \(b\), up to
  local contact terms already fixed by the Ward identity.  The free Dirac
  determinant gives the transverse quadratic functional
  \begin{equation}
    W_E[b]
    =
    \frac1{2\pi}
    \int\frac{\dd^2p}{(2\pi)^2}\,
    b_\mu(-p)
    \left(\delta^{\mu\nu}-\frac{p^\mu p^\nu}{p^2}\right)
    b_\nu(p).
    \label{eq:schwinger-euclidean-current-functional}
  \end{equation}
  The same expression is obtained by integrating out the Euclidean scalar in
  the right side of \eqref{eq:schwinger-current-generating-functional},
  because
  \[
    \epsilon^{\mu\alpha}p_\alpha
    \epsilon^{\nu\beta}p_\beta
    =
    p^2\delta^{\mu\nu}-p^\mu p^\nu .
  \]
  The Abelian current algebra of a free Dirac fermion is Gaussian: after the
  current-current short-distance singularity is fixed, higher connected
  current correlators have no independent singular tensor structure
  compatible with current conservation and chirality.  Thus
  \eqref{eq:schwinger-euclidean-current-functional} determines the current
  sector.  Analytic continuation gives
  \eqref{eq:schwinger-current-generating-functional}.
\end{proof}

The scalar in this theorem is compact.  The derivative operators
\(\partial_\mu\varphi\) describe the current algebra, while charged fermion
operators require vertex operators.  With the present normalization the
fermion mass operator is proportional to
\(\cos(2\sqrt\pi\,\varphi)\); therefore
\(\varphi\sim\varphi+\sqrt\pi\).  This compactness is invisible if one looks
only at the local current \(\partial_\mu\varphi\), but it is essential for
the charge lattice and for the massive deformation.

\section{Exact Bosonized Gauge Theory}

Setting \(b_\mu=e a_\mu\) in
\eqref{eq:schwinger-current-generating-functional} and adding the Maxwell
term and the theta term gives the exact bosonized representative
\begin{equation}
  \mathcal L_{\rm bos}
  =
  -\frac12\partial_\mu\varphi\,\partial^\mu\varphi
  -\frac14F_{\mu\nu}F^{\mu\nu}
  +\frac{e}{\sqrt\pi}a_\mu\epsilon^{\mu\nu}\partial_\nu\varphi
  +\frac{e\theta}{2\pi}E .
  \label{eq:schwinger-bosonized-action}
\end{equation}
This is an exact equality for the local current and electric-field sector of
the massless theory.  It should not be read as a statement that a
gauge-variant fermion field has become a gauge-invariant local scalar
operator.  The charged sectors are encoded by the compact-boson vertex
operators together with the gauge line required to make a charged insertion
gauge invariant.

Up to a boundary term,
\[
  a_\mu\epsilon^{\mu\nu}\partial_\nu\varphi=E\varphi .
\]
Therefore
\begin{equation}
  \mathcal L_{\rm bos}
  =
  -\frac12\partial_\mu\varphi\,\partial^\mu\varphi
  +\frac12E^2
  +E\left(\frac{e}{\sqrt\pi}\varphi+\frac{e\theta}{2\pi}\right).
  \label{eq:schwinger-electric-field-form}
\end{equation}
The field \(E\) has no independent transverse wave degree of freedom in two
spacetime dimensions.  The stationary equation for \(E\) on a fixed local
branch is
\[
  E
  =
  -\frac{e}{\sqrt\pi}\varphi-\frac{e\theta}{2\pi}.
\]
Eliminating \(E\) gives
\begin{equation}
  \mathcal L_{\rm eff}
  =
  -\frac12\partial_\mu\varphi\,\partial^\mu\varphi
  -\frac12\,\frac{e^2}{\pi}
  \left(\varphi+\frac{\theta}{2\sqrt\pi}\right)^2 .
  \label{eq:schwinger-effective-massive-scalar}
\end{equation}
The exact gauge-invariant local spectrum therefore contains one neutral
scalar particle of mass
\begin{equation}
  m_{\rm Sch}^2=\frac{e^2}{\pi}.
  \label{eq:schwinger-mass}
\end{equation}
The electric field is the corresponding massive field,
\begin{equation}
  E
  =
  -\frac{e}{\sqrt\pi}
  \left(\varphi+\frac{\theta}{2\sqrt\pi}\right),
  \qquad
  (\Box-m_{\rm Sch}^2)E=0 .
  \label{eq:schwinger-electric-field-scalar-relation}
\end{equation}

\begin{remark}[Theta periodicity and branches]
  The local expression \eqref{eq:schwinger-effective-massive-scalar} is
  written on one electric-flux branch.  In the compact theory, large gauge
  transformations and the compactness
  \(\varphi\sim\varphi+\sqrt\pi\) identify the branch obtained by
  \(\theta\mapsto\theta+2\pi\) with the branch obtained by shifting
  \(\varphi\).  A formulation on a spatial circle keeps the holonomy zero
  mode and sums over the corresponding integer electric-flux sectors.  The
  massless theory is independent of \(\theta\), because an axial rotation
  shifts \(\theta\) and no fermion mass term obstructs this shift.
\end{remark}

\begin{figure}[H]
  \centering
  \resizebox{0.96\textwidth}{!}{%
  \begin{tikzpicture}[>=Stealth,
    box/.style={draw,rounded corners=2pt,align=center,inner sep=5pt,
      text width=3.35cm,minimum height=1.05cm},
    arrowlabel/.style={font=\scriptsize,fill=white,inner sep=1pt}]
    \node[box] (ferm) at (-5.2,0)
      {massless Dirac field\\and Abelian gauge representative};
    \node[box,text width=3.75cm] (curr) at (0,0)
      {current algebra\\\(J^\mu=\pi^{-1/2}\epsilon^{\mu\nu}\partial_\nu\varphi\)};
    \node[box] (mass) at (5.2,0)
      {one massive scalar\\\(m_{\rm Sch}^2=e^2/\pi\)};
    \draw[->,thick] (ferm.east) -- node[arrowlabel,above] {bosonization} (curr.west);
    \draw[->,thick] (curr.east) -- node[arrowlabel,above] {integrate \(E\)} (mass.west);
    \node[align=center] at (0,-1.15)
      {\(\partial_\mu J_A^\mu=(e/\pi)E\) fixes the mass term};
  \end{tikzpicture}
  }
  \caption{The exact solution of the local sector of the massless Schwinger
  model.  The anomaly and current algebra convert the gauge-invariant
  electric field into a free massive scalar field.}
  \label{fig:schwinger-model-exact-solution}
\end{figure}

\section{Static Probes and Screening}

The response to external static charges is a line-observable question, not a
statement about the local particle spectrum alone.  For a neutral pair
\[
  \rho_{\rm ext}(x)
  =
  Q\delta(x-R/2)-Q\delta(x+R/2),
\]
pure Maxwell theory in one spatial dimension gives a constant electric field
between the probes.  With the Gauss law \(\partial_xE+\rho_{\rm ext}=0\) in
the interval between dynamical charges,
\[
  V_{\rm pure}(R)=\frac{Q^2}{2}R .
\]

The massless Schwinger model changes the static kernel.  Integrating out the
massless Dirac current sector gives, for the static temporal representative
in Coulomb gauge, the quadratic energy functional
\begin{equation}
  \mathcal E[a_0;\rho_{\rm ext}]
  =
  \frac12\int\dd x\,
  a_0(x)\left(-\partial_x^2+m_{\rm Sch}^2\right)a_0(x)
  -\int\dd x\,a_0(x)\rho_{\rm ext}(x).
  \label{eq:schwinger-static-energy-functional}
\end{equation}
The Euler-Lagrange equation is
\[
  \left(-\partial_x^2+m_{\rm Sch}^2\right)a_0(x)
  =
  \rho_{\rm ext}(x).
\]
Let
\[
  \left(-\partial_x^2+m^2\right)G_m(x)=\delta(x),
  \qquad
  G_m(x)=\frac{\ee^{-m|x|}}{2m}.
\]
Solving for \(a_0\), substituting back into
\eqref{eq:schwinger-static-energy-functional}, and subtracting the
\(R\)-independent self-energies gives
\begin{equation}
  V_{\rm Sch}(R)
  =
  Q^2\left(G_{m_{\rm Sch}}(0)-G_{m_{\rm Sch}}(R)\right)
  =
  \frac{Q^2}{2m_{\rm Sch}}
  \left(1-\ee^{-m_{\rm Sch}R}\right).
  \label{eq:schwinger-screened-static-potential}
\end{equation}
Thus the force decays exponentially and the energy approaches a constant at
large separation.  The constant is the energy needed to polarize the
massless charged vacuum around the two probes.  There is no asymptotic
linear electric string in the massless model.

This screening statement is compatible with
\eqref{eq:schwinger-electric-field-mass-equation}.  The local electric field
creates a massive neutral particle, while external flux is screened because
the induced charge density can rearrange over the length scale
\(m_{\rm Sch}^{-1}\).  These are two different gauge-invariant observables
whose common origin is the same anomaly-fixed current algebra.

\section{Massive Fermion Variant and Fractional-Probe Confinement}

Now add a Dirac mass \(M\bar\psi\psi\).  The axial rotation that removed
\(\theta\) in the massless theory also rotates the mass phase, so the
physical parameter is the relative phase between the mass and the theta
term.  Choosing \(M>0\) real, bosonization gives the sine-Gordon perturbation
\begin{equation}
  \Delta\mathcal L_M
  =
  \kappa M\cos(2\sqrt\pi\,\varphi+\theta).
  \label{eq:massive-schwinger-sine-gordon-perturbation}
\end{equation}
The constant \(\kappa>0\) has mass dimension one and is fixed by the
normalization of the renormalized composite mass operator.  Equivalently,
\(\kappa\) is the magnitude of the chiral condensate in the massless theory
in the corresponding normal-ordering scheme.  The scheme dependence of
\(\kappa\) is exactly compensated by the scheme dependence of the
renormalized mass \(M\).

After eliminating \(E\) on a local branch, the scalar potential is
\begin{equation}
  V_M(\varphi)
  =
  \frac12\,\frac{e^2}{\pi}
  \left(\varphi+\frac{\theta}{2\sqrt\pi}\right)^2
  -\kappa M\cos(2\sqrt\pi\,\varphi+\theta).
  \label{eq:massive-schwinger-effective-potential}
\end{equation}
The quadratic term should again be understood together with the integer flux
branches of the compact gauge theory.  The cosine is periodic under
\(\varphi\mapsto\varphi+\sqrt\pi\); the electric term records which flux
branch is being used.

External charges shift the effective theta angle in the interval between
them.  This is the cleanest derivation of the fractional-probe string
tension.  A charge \(Q\) at one endpoint and \(-Q\) at the other force the
electric flux in the interval to differ from the outside flux by \(Q\).  In
the bosonized description this is equivalent to
\begin{equation}
  \theta_{\rm inside}
  =
  \theta+\frac{2\pi Q}{e}
  \quad
  \operatorname{mod} 2\pi .
  \label{eq:schwinger-theta-shift-probe}
\end{equation}
For small \(M\), the leading vacuum energy density of the massive theory is
computed by evaluating the mass perturbation in the massless vacuum:
\[
  \varepsilon(\theta)
  =
  -\kappa M\cos\theta+O(M^2).
\]
Therefore the fixed-branch large-\(R\) energy-density shift in the interval,
relative to the outside vacuum, is
\begin{equation}
  \Delta\varepsilon_Q(\theta)
  =
  \varepsilon\!\left(\theta+\frac{2\pi Q}{e}\right)-\varepsilon(\theta)
  =
  \kappa M
  \left[
    \cos\theta-\cos\!\left(\theta+\frac{2\pi Q}{e}\right)
  \right]
  +O(M^2),
  \label{eq:massive-schwinger-string-tension-theta}
\end{equation}
with the physical branch chosen by minimizing over the integer flux sectors.
At \(\theta=0\) this fixed-branch shift is nonnegative and gives the leading
fractional-probe string tension
\begin{equation}
  \sigma(Q)
  =
  \kappa M
  \left(1-\cos\frac{2\pi Q}{e}\right)
  +O(M^2).
  \label{eq:massive-schwinger-string-tension}
\end{equation}
For \(Q/e\in\mathbb Z\), the string tension vanishes: an integer external
charge lies in the dynamical charge lattice and can be screened.  For
nonintegral \(Q/e\), the interval between the probes is forced into a
different theta branch, and its energy grows linearly at large separation
with leading coefficient \(\sigma(Q)\).  This is confinement of fractional
test charges, not confinement of the dynamical unit-charged fermion.

\section{Mechanisms Isolated by the Exact Solution}

The model is useful because several mechanisms that are intertwined in
higher-dimensional gauge theories become separately computable:
\[
\begin{array}{c|c}
  \hbox{object} & \hbox{exact statement in the Schwinger model} \\ \hline
  \hbox{local anomaly} &
  \partial_\mu J_A^\mu=(e/\pi)E \\
  \hbox{local electric-field spectrum} &
  E\hbox{ creates a scalar of }m_{\rm Sch}^2=e^2/\pi \\
  \hbox{massless probe response} &
  V(R)=Q^2(1-\ee^{-m_{\rm Sch}R})/(2m_{\rm Sch}) \\
  \hbox{mass perturbation} &
  M\bar\psi\psi\leftrightarrow\kappa M\cos(2\sqrt\pi\varphi+\theta) \\
  \hbox{fractional-probe confinement} &
  \sigma(Q)=\kappa M(1-\cos(2\pi Q/e))+O(M^2)
\end{array}
\]
Each line refers to a specified gauge-invariant local observable or to an
explicit external line probe.  This separation is essential: a massive
electric-field two-point function, a screened Wilson-line potential, and a
linear potential for fractional probes in the massive deformation are
different physical statements, even though in this exactly solvable theory
they are controlled by the same current algebra and anomaly.
